\documentclass[10pt,twocolumn,letterpaper]{article}

\usepackage{cvpr}
\usepackage{times}
\usepackage{epsfig}
\usepackage{graphicx}
\usepackage{amsmath}
\usepackage{amssymb}

% Include other packages here, before hyperref.

% If you comment hyperref and then uncomment it, you should delete
% egpaper.aux before re-running latex.  (Or just hit 'q' on the first latex
% run, let it finish, and you should be clear).
\usepackage[pagebackref=true,breaklinks=true,letterpaper=true,colorlinks,bookmarks=false]{hyperref}

\cvprfinalcopy % *** Uncomment this line for the final submission

\def\cvprPaperID{42} % *** Enter the CVPR Paper ID here
\def\httilde{\mbox{\tt\raisebox{-.5ex}{\symbol{126}}}}

% Pages are numbered in submission mode, and unnumbered in camera-ready
\ifcvprfinal\pagestyle{empty}\fi
\begin{document}

%%%%%%%%% TITLE
\title{Scene Classification with Deep Convolutional Neural Networks}

\author{Yangzihao Wang and Yuduo Wu\\
University of California, Davis\\
{\tt\small \{yzhwang,yudwu\}@ucdavis.edu}
% For a paper whose authors are all at the same institution,
% omit the following lines up until the closing ``}''.
% Additional authors and addresses can be added with ``\and'',
% just like the second author.
% To save space, use either the email address or home page, not both
}
\maketitle
%\thispagestyle{empty}

%%%%%%%%% ABSTRACT
\begin{abstract}
%summarize the problem, main idea, and results;
The use of massive datasets like ImageNet and the revival of Convolutional
Neural Networks (CNNs) for learning deep features has significantly improved
the performance of object recognition. However, performance at scene
classification has not achieved the same level of success since there is still
semantic gap between the deep features and the high-level context.  In this
project we proposed a novel scene classification method which combines CNN and
Spatial Pyramid to generate high-level context-aware features for one-vs-all
linear SVMs. Our method achieves the state-of-the-art result: 68.04\% average
accuracy rate on MIT indoor67 dataset using only the deep features trained from
ImageNet.

\end{abstract}

%%%%%%%%% BODY TEXT
\section{Related Work}
\label{sec:related}
%provide a detailed description of related papers (not necessarily limited to
%those in the schedule).  If you're proposing a new idea or extending an existing
%approach, compare and contrast it with existing work.  If you're analyzing one
%or two related techniques, describe how they relate to other relevant work;

Scene classification means to provide information about the semantic category
or the function of a given image. Among different kidn of scene classification
tasks, the indoor scene classification is considered to be one of the most
difficult since the lack of discriminative features and contexts at the high
level~\cite{Quattoni:2009:RIS}. Spatial pyramid
representation\cite{Lazebnik:2006:BBF} is a popular method used for scene
classification tasks. It is a simple and computationally efficient extension of
an orderless bag-of-features image representation. However, without a proper
high-level feature representation, such schemes often fail to offer sufficient
semantic information of a scene. Object bank\cite{Li:2010:OBA} is among the
first to propose a high-level image representation for scene classification. It
uses a large number of pre-trained generic object detectors to create response
maps for high level visual recognition tasks. The combination of off-the-shelf
object detectors and a simple linear prediction model with a sparse-coding
scheme achieves superior predictive power over similar linear prediction models
trained on conventional representations. However, this method also limits the
performance of their system to the performance of the object detectors they
choose. Recently, Convolutional Neural Networks (CNNs) with flexible capacity
makes training from large-scale dataset such as ImageNet~\cite{Deng:2009:IAL}
possible. In the work of A. Krizhevsky et al.\cite{Krizhevsky:2012:ICD}, they
trained one of the largest CNNs on the subsets of ImageNet and achieved better
results than any other state-of-the-art methods in 2012. While their CNN system
focuses on object detection, the features generated can be used for other
applications such as scene classification. Two types of improvements has been
done on top of their CNN works. The first type of improvement tries to address
the problem of generating possible object locations in an image. Selective
search method~\cite{Uijlings:2013:SSO} combines the strength of both an
exhaustive search and segmentation and results in a small set of data-driven,
class-independent, high quality locations. Girshick et al. propose the Regions
with CNN features (R-CNN) method~\cite{Girshick:2013:RFH} as a more effective
feature generation method. Alternatively, Zhou et al. try to increase the
performance of scene classification using CNN by creating a new scene-centric
database~\cite{Zhou:2014:LDF}.


%%%%%%%%% BODY TEXT
\section{Technical Approach}
\label{sec:method}
%Describe in detail the feature representation(s) and algorithm(s) you employed.
%The description should be self-contained (i.e., the reader should not have to
%rely on outside sources for your points to be clear), and should provide enough
%detail so that the reader could re-implement the approach. Clearly state the
%method's input and output, and any assumptions or design choices;

Previous work on Convolutional Neural Networks (CNNs) implies that it may capture the high-level representation of an image
using a certain deep layer feature set. Our goal of this project is to answer one single question: \emph{Whether CNNs can help with the feature representation
to extract high-level inforamtion of an image scene and thus improve the scene classification precision?} 
We choose a CNN which is pre-trained on ImageNet dataset (ImageNet-CNN) since its a large-scale general object recognition datasets which consists of over 15 million labeled
high-resolution images in over 22,000 categories. We use CNN pretrained on such dataset with the hope to reduce the chance of overfitting to certain scenes.
To utilize a pre-trained ImageNet CNN and for the efficiency of the feature extraction process, we use a popular library: Caffe~\cite{Jia:2014:CCA}.

For the training process, our system takes all images in the training set for each category as the input, use the ImageNet-CNN to perform a prediction for each image. Instead of
getting the final 1000 length prediction vector, we extract the FC 7 layer feature set which contains 4096 response values. We then use such features to train one linear SVM model
for each scene category. For the testing process, an input image goes through the same ImageNet-CNN and its 4096 length deep feature vector are used to predict its scene classification
for each linear SVM model and we assign the one with highest confidence score.

Within this general framework, several methods can be explored to improve the feature representation. Instead of using the entire image for deep feature extraction, we can first
select a set of region proposals (usually around 2000 for a high-resolution image) which are most informative about the image, then extract 4096-dimensional feature vectors for each region proposal. This improvement puts more weights on more informative regions. However, using the concadenated feature vector would result in an extremely high dimension (2000 $\times$ 4096), it is necessary to use some sparse coding scheme to represent the global feature pattern. In this paper, we adopt spatial pyramid because of its simplicity and effectiveness. To increase the generalization ability of our features and to reduce the impact of overfitting, we further apply $l-2$ normalization to the achieved feature vectors before we feed them to the one-vs-all SVMs. Figure~\ref{fig:system_overview} shows the pipeline of our system. We will describe the details in the following sections.

\begin{figure*}[ht]
  \centering
  \includegraphics[scale=0.8]{img/demo.pdf}
  \centering
  \caption{The overview of our system. For an input image, a selective search algorithm is applied first to get roughly 2000 regions of interest. We then apply a pre-trained Convolutional Neural Network (CNN) on each region of interest to get a deep feature vector of length 4096. A three-level spatial pyramid representation of the image with deep feature is used to create the final feature representation. At each level, for each spatial bin, we use max pooling to get the largest feature value of all the feature values of the regions of interest which fall into that spatial bin, resulting in the final feature of length No.Deep Features $\times$ (1+4+16) as a high-level representation of the input image. Then multiple one-vs-all linear SVMs are used to do the scene classification.}
  \label{fig:system_overview}
\end{figure*}

\subsection{Selective Search}
A variety of recent research offers methods for generating category-independent
region proposals for possible object locations.
Selective search is widely used for generating possible object locations for
use in object recognition\cite{Uijlings:2013:SSO}. Same strategies can be
adopted on indoor scene classification. For the indoor scenes that can be well
characterized by objects they contain, the selective search can exploit local
discriminative information with greatly reduced number of locations compared
to an exhaustive search. We use selective search to generate region proposals.
Caffe provides a general Python interface for models and it has built in
interface for selective search. We only need to change the setting of CROP\_MODES
to selective\_search, we can operate on around 2000 region proposals instead of
the entire image.

\subsection{Feature Extraction}
Caffe~\cite{Jia:2014:CCA} is an open source convolutional architecture for
fast feature embedding which contains pre-trained models.
In our project, we use pre-trained BVLC Reference CaffeNet to classify images
and extract 4096-dimensional layer 7 feature vectors from each region
proposal using Caffe~\cite{Jia:2014:CCA} implementation of the CNN described by
Krizhevsky et al\cite{Krizhevsky:2012:ICD}. It provides an option to output the
features in certain layer rather than only the final classification results. We
set the \textit{blobs} option to \textit{fc7} in order to obtained the Layer
7 feature vectors for each of the region proposals. Therefore, after this step,
for one input image, we obtain around 2000 feature vectors and each have the
dimension of 4096.

\subsection{Max Pooling}

\subsection{Spatial Pyramid Matching}
For each image, a three-level spatial pyramid representation is used, resulting
$numImages * numWindows * (1 + 4 + 16)$ length feature vectors.

\subsection{L2 Normalization}
We then perform the L2 normalizations for each of the feature matrix.
L2 normalization is computed by the square root of the sum of each feature's
square and for each feature in the feature vector. By dividing each feature
vector the L2 normalization value, we have each feature vector's
L2 normalization = 1. This L2 normalizations have only a modest effect on the
overall performance perhaps simply because the original feature vector data
already in the similar scale.

\subsection{Training}
SVMs (Support Vector Machines) are a useful technique for data classification.
LibSVM\cite{Chang:2011:CC01a} is an integrated software for support vector
classification that is widely used in variety of classification tasks.
It supports multi-class classification which is used in this project.
We trained 67 binary one-vs-all SVM classifiers each for one category in
MIT-indoor67 dataset.
In order to fit the required LibSVM training file format, for each category
training file, we add label +1 to each feature vector that belongs to the
category and label -1 to all feature vectors that belong to rest of categories.
We also move all instances that belonging to the current positive category
(+1 labeled feature vectors) to be at the top of the feature matrix, this would
guarantee the correctness even if LibSVM might internally map the label of the
first training instance to be +1 regardless of its actual label value.


\section{Experiments}
\label{sec:results}
Describe the experiments you conducted to evaluate the approach.  For each
%experiment, describe what you did, what was the main purpose of the experiment,
%and what you learned from the results. Provide figures, tables, and qualitative
%examples, as appropriate.

In this section, we evaluate our method on the indoor scene dataset:
MIT-indoor67~\cite{Quattoni:2009:RIS},
which includes 15,620 images over 67 indoor scenes. Suggested training and
testing list of images are used to do the training (80 images per class)
and validation (20 images per class). There are at least 100 images per
scene and all in .jpg format.

Multi-class classification is done with a support vector machine (SVM) trained
using one-versus-all rule, that is, each classifier is learned to separate each
class from the rest of classes. Test image is assigned the label of the
classifier with the highest response.Scene classification performance is
evaluated by average multi-class classification accuracy over all scene classes.

For comparison purposes, we implement with the same procedure but only use
the extracted layer 7 4096-dimensional feature vectors from Caffe. After we get
one feature vectors for each entire image, instead of perform spatial pyramid
and L2 normalization, we simply add labels and send them into the multi-class
SVMs. Validation image feature vectors are also generated in the same way.

Our
method achieves a mean average precision (mAP) of 68.2953\% on dataset
MIT-indoor67\cite{Quattoni:2009:RIS} without fine-tune on this dataset.
For comparison, we implement using only 4096-dimensional feature vectors
extracted from Caffe without region proposals, spatial pyramid matching 
and max-pooling which has the mAP of 59.9507\%.

% Table 1
\begin{table}[ht]
	\caption{Comparison results on MIT-indoor67}
	\centering
	\begin{tabular}{l c c}
	\hline \hline
	Models                & Average Precision \\ \hline
	Our Method  	      & {\bf{68.2953\%}} \\
	Without Normalization & 68.0469\% \\
        Only CNN Feature      & 59.9507\% \\
	\hline
	\end{tabular}
	\label{tab:overall}
\end{table}

% overall results
We compare our scene classification tasks with the performance without perform
L2 normalization and the performance of using only the features extracted from
CNN, summarized in Table~\ref{tab:overall}. Around 14\% improvement of using
region proposals, max-pooling and spatial pyramid are shown. In majority of
categories, we perform much better on the average precision. Some examples
are shoes shop, bedroom, bowling, grocery store, hospital room and operating
room. This might due to the region proposals and spatial pyramid technique
allow us to better characterize the particular objects belong to the category.
However, there are also some drops of average accuracy using our methods and
mainly for these three categories: prison-cell, library and living room. These
three categories are all relatively easier to be characterized by global
spatial properties.

\iffalse
% Table 2
\begin{table}[ht]
        \caption{Average Acurracy that Drops}
        \centering
        \begin{tabular}{l c c}
        \hline \hline
        Category    & our method & only feature \\ \hline
        prison-cell & 45\%       & 70\% \\
        library     & 45\%       & 70\% \\
        livingroom  & 20\%       & 40\% \\
        \hline
        \end{tabular}
        \label{tab:overall}
\end{table}
\fi
% compare with other paper results

% heatmap visulizations


%-------------------------------------------------------------------------
\section{Conclusions}

briefly summarize the main idea and results, and possible future work.

{\small
\bibliographystyle{ieee}
\bibliography{report}
}

\appendix
\section*{Appendix: Confusion Matrix}

\label{sec:appendix}
% confusion matrix
\FloatBarrier
\begin{figure}[h!]
  \centering
  \includegraphics[width=0.95\textwidth]{img/matrix.pdf}
  \centering
  \caption{Confusion matrix of prediction results for 67 categories.}
  \label{fig:confusion_matrix}
\end{figure}
\FloatBarrier

\end{document}